\section*{Sets and Indices}

\begin{itemize}
\item $I$ : applicant types, indexed by $i$.
\[
I=\{1,2,3,4,5,6\}.
\]

\item $J$ : branch cities, indexed by $j$.
\[
J=\{\text{Donghai},\text{Nanjiang}\}.
\]

\item $K$ : specialties, indexed by $k$.
\[
K=\{1,2,3\}.
\]

\item $S_i \subseteq K$ : specialties suitable for applicant type $i$ (code: \texttt{supply[i]['suit']}).
\end{itemize}

\section*{Parameters}

\begin{itemize}
\item $N_i$ : number of available applicants of type $i$ (code: \texttt{supply[i]['num']}).

From Table 4.4,
\[
N_i=1500 \qquad \forall i\in I.
\]

\item $D_{jk}$ : demand for specialty $k$ in city $j$ (code: \texttt{demand[(j,k)]}).

From Table 4.3,
\[
\begin{aligned}
&D_{\text{Donghai},1}=1000,\quad D_{\text{Donghai},2}=2000,\quad D_{\text{Donghai},3}=1500,\\
&D_{\text{Nanjiang},1}=2000,\quad D_{\text{Nanjiang},2}=1000,\quad D_{\text{Nanjiang},3}=1000.
\end{aligned}
\]

\item $a_i$ : preferred specialty of applicant type $i$ (code: \texttt{supply[i]['pref\_spec']}).

\item $b_i$ : preferred city of applicant type $i$ (code: \texttt{supply[i]['pref\_city']}).

\item $T_2$ : target number meeting preferred specialty (code: \texttt{p2\_target}).
\[
T_2=8000.
\]

\item $T_3$ : target number meeting preferred city (code: \texttt{p3\_target}).
\[
T_3=8000.
\]

\item $\delta_{ik}$ : suitability indicator, equal to $1$ if $k\in S_i$, and $0$ otherwise.
This is implemented through zero-fixing constraints rather than explicit coefficients.
\end{itemize}

\section*{Decision Variables}

For each solve, the model uses integer assignment variables and deviational variables.

\begin{itemize}
\item $x_{ijk}$ (code: \texttt{x1[i,j,k]} in Phase 1, \texttt{x2[i,j,k]} in Phase 2): number of applicants of type $i$ assigned to city $j$ and specialty $k$.
\[
x_{ijk}\in \mathbb{Z}_{\ge 0}.
\]

\item $d_2^-$, $d_2^+$ (code: \texttt{d2\_minus1}, \texttt{d2\_plus1} in Phase 1; \texttt{d2\_minus2}, \texttt{d2\_plus2} in Phase 2): negative and positive deviations for Priority 2.
\[
d_2^-,d_2^+\in \mathbb{Z}_{\ge 0}.
\]

\item $d_3^-$, $d_3^+$ (code: \texttt{d3\_minus2}, \texttt{d3\_plus2} in Phase 2): negative and positive deviations for Priority 3.
\[
d_3^-,d_3^+\in \mathbb{Z}_{\ge 0}.
\]
\end{itemize}

\section*{Objective}

The code implements a two-phase lexicographic goal program.

\paragraph{Phase 1.}
Minimize the shortfall from the Priority 2 target:
\[
\min d_2^- .
\]

\paragraph{Phase 2.}
Let $z_2^\star$ be the optimal value of $d_2^-$ from Phase 1 (code: \texttt{min\_d2}). Then minimize the shortfall from the Priority 3 target subject to not worsening Phase 1:
\[
\min d_3^- .
\]

\section*{Constraints}

\paragraph{Phase 1 model.}
\[
\min \ d_2^-
\]
subject to
\[
\sum_{j\in J}\sum_{k\in K} x_{ijk} \le N_i \qquad \forall i\in I,
\]
\[
\sum_{i\in I} x_{ijk} \ge D_{jk} \qquad \forall j\in J,\ \forall k\in K,
\]
\[
x_{ijk}=0 \qquad \forall i\in I,\ \forall j\in J,\ \forall k\in K\setminus S_i,
\]
\[
\sum_{i\in I}\sum_{j\in J} x_{ij,a_i} + d_2^- - d_2^+ = T_2,
\]
\[
x_{ijk}\in \mathbb{Z}_{\ge 0},\qquad d_2^-,d_2^+\in \mathbb{Z}_{\ge 0}.
\]

\paragraph{Phase 2 model.}
\[
\min \ d_3^-
\]
subject to
\[
\sum_{j\in J}\sum_{k\in K} x_{ijk} \le N_i \qquad \forall i\in I,
\]
\[
\sum_{i\in I} x_{ijk} \ge D_{jk} \qquad \forall j\in J,\ \forall k\in K,
\]
\[
x_{ijk}=0 \qquad \forall i\in I,\ \forall j\in J,\ \forall k\in K\setminus S_i,
\]
\[
\sum_{i\in I}\sum_{j\in J} x_{ij,a_i} + d_2^- - d_2^+ = T_2,
\]
\[
d_2^- \le z_2^\star,
\]
\[
\sum_{i\in I}\sum_{k\in K} x_{i,b_i,k} + d_3^- - d_3^+ = T_3,
\]
\[
x_{ijk}\in \mathbb{Z}_{\ge 0},\qquad d_2^-,d_2^+,d_3^-,d_3^+\in \mathbb{Z}_{\ge 0}.
\]

Here,
\[
\sum_{i\in I}\sum_{j\in J} x_{ij,a_i}
\]
counts recruits assigned to their preferred specialty, and
\[
\sum_{i\in I}\sum_{k\in K} x_{i,b_i,k}
\]
counts recruits assigned to their preferred city.

\section*{Notes on Implementation}

\begin{itemize}
\item The implementation does \emph{not} encode Priority 1 as an objective. Instead, it enforces branch-specialty requirements by the hard constraints
\[
\sum_i x_{ijk} \ge D_{jk}.
\]
Thus all stated demands must be satisfied in every feasible solution.

\item Demand constraints are inequalities ``$\ge$'', not equalities. Hence the model permits over-assignment to a city-specialty pair if useful for later priorities.

\item Supply constraints are inequalities ``$\le$'', so some applicants may remain unassigned.

\item Suitability is enforced exactly as in code by fixing $x_{ijk}=0$ whenever specialty $k$ is not in the suitable-specialty list of type $i$.

\item The goal-programming logic is sequential. Phase 1 minimizes only the negative deviation for Priority 2. Phase 2 then minimizes only the negative deviation for Priority 3 while imposing
\[
d_2^- \le z_2^\star.
\]
Because Phase 1 already minimizes $d_2^-$, this reproduces the implemented lexicographic priority order.

\item The code does not penalize positive deviations $d_2^+$ or $d_3^+$. Therefore exceeding either target is allowed at no direct objective cost.

\item All variables are declared integer and the models are solved as MIPs with Gurobi through PuLP compatibility. The reported answer is the Phase 2 optimum $d_3^-$, interpreted as the minimum number of recruited people that cannot meet Priority 3 while respecting the earlier priority structure implemented in the code.
\end{itemize}
